\documentclass[11pt]{article}

\usepackage[margin=1in]{geometry}
\usepackage{amsmath,amssymb,booktabs,tabularx}

\title{Riccati/ODE system for \((C,\mathbf D,V)\)}
\author{}
\date{}

\begin{document}
\maketitle

\noindent\textbf{Status.} Draft\\
\textbf{Plan item.} Riccati/ODE system for \((C,D,V)\)\\
\textbf{Depends on.} \texttt{04\_affine\_transform\_ansatz.tex}\\
\textbf{Feeds.} Affine schedule transform, affine domain constraints, terminal pricing measure

\section{Purpose}

Derive the deterministic equations for the affine schedule loadings
\((C,\mathbf D,V)\). These equations are the actual transform recursion used
once the affine ansatz has reduced the conditional expectation to a finite
system of coefficient equations.

\section{Setup}

Use the same one-interval setup as
\texttt{04\_affine\_transform\_ansatz.tex}. Let \(\tau\in[0,\Delta]\) denote
backward elapsed time from the terminal condition. The terminal conditions are
\[
C(0)=0,
\qquad
\mathbf D(0)=\mathbf D_{\mathrm{term}},
\qquad
V(0)=\xi.
\]

The Hawkes jump vectors and news loading vector are
\[
\mathbf{k}_+=(\alpha_{++},0,\alpha_{-+},0),
\qquad
\mathbf{k}_-=(0,\alpha_{+-},0,\alpha_{--}),
\qquad
\mathbf r=(\rho_{++},\rho_{+-},\rho_{-+},\rho_{--}).
\]

\section{Notation}

\begin{tabularx}{\textwidth}{lXX}
\toprule
Symbol & Meaning & Units or domain\\
\midrule
\(H_\pm^B(\mathbf D,V)\) & Baseline-event Hamiltonian after baseline clock loading is included & Complex finite\\
\(H_\pm^E(\mathbf D,V)\) & Endogenous-event Hamiltonian after endogenous clock loading is included & Complex finite\\
\(\kappa_z\) & Long-variance mean-reversion rate & Strictly positive\\
\(\bar z\) & Derived long-variance mean-reversion level & Strictly positive\\
\(\sigma_z\) & Long-variance volatility of volatility & Non-negative\\
\(\chi_\pm\) & Long-variance event-jump loading by side & Non-negative\\
\(M_c(q,r)\) & Gamma-normal news mark transform for component \(c\) & Moment domain\\
\(\ell_c(q,r)\) & \(\log M_c(q,r)\) on the principal branch & Moment domain\\
\(\nu_k^{\mathrm{known}}\) & Equivalent rate for scheduled known-news timestamps & Non-negative\\
\(\nu_k^{\mathrm{surprise}}\) & Surprise-news Poisson rate & Non-negative\\
\bottomrule
\end{tabularx}

\section{Assumptions used}

\begin{tabularx}{\textwidth}{XXX}
\toprule
Assumption & Where it is used & Why it is needed\\
\midrule
The affine ansatz holds for the interval transform.
& Coefficient matching between \(\partial_\tau f/f\) and \(\mathcal A f/f\).
& This turns the transform problem into ODEs for \(C\), \(\mathbf D\), and \(V\).\\
Buy endogenous intensity is \(Y^{++}+Y^{+-}\), and sell endogenous intensity is \(Y^{-+}+Y^{--}\).
& Assigning \(H_+^E\) to \(D_{++},D_{+-}\) and \(H_-^E\) to \(D_{-+},D_{--}\).
& The side intensity convention determines which state coordinates get each Hamiltonian.\\
Baseline event intensities are both \(s_k\bar\lambda\).
& Deriving the scalar baseline event term and the \(Z\)-linear baseline clock term.
& Baseline arrivals affect \(C\), while baseline return clocks affect \(V\).\\
The long-variance diffusion is CIR affine.
& Deriving the Riccati terms in \(V'\) and the drift term in \(C'\).
& A square-root affine diffusion produces \(V^2\) but no higher powers of \(Z\).\\
Known news is deterministic-arrival random-mark and surprise news is compound Poisson.
& Choosing \(\nu^{\mathrm{known}}\ell_{\mathrm{known}}\) and \(\nu^{\mathrm{surprise}}(M_{\mathrm{surprise}}-1)\).
& The two news components have different arrival transforms.\\
\bottomrule
\end{tabularx}

\section{Result}

\noindent\textbf{Identity.} Define
\[
\begin{aligned}
H_+^B(\mathbf D,V)
&=
z_+\exp(\mathbf D\cdot\mathbf k_+ + V\chi_+) - 1 - u_+^B,\\
H_-^B(\mathbf D,V)
&=
z_-\exp(\mathbf D\cdot\mathbf k_- + V\chi_-) - 1 - u_-^B,\\
H_+^E(\mathbf D,V)
&=
z_+\exp(\mathbf D\cdot\mathbf k_+ + V\chi_+) - 1 - u_+^E,\\
H_-^E(\mathbf D,V)
&=
z_-\exp(\mathbf D\cdot\mathbf k_- + V\chi_-) - 1 - u_-^E.
\end{aligned}
\]
For \(c\in\{\mathrm{known},\mathrm{surprise}\}\), define
\[
M_c(q,r)
=
\left(
1-\theta_U\left(ra_c+qm_c+\frac{1}{2}q^2\sigma_c^2\right)
\right)^{-k_U},
\qquad
\ell_c(q,r)=\log M_c(q,r).
\]
Then the affine schedule loadings solve
\[
\begin{aligned}
D_{++}'(\tau) &= -\beta D_{++}(\tau)
    +H_+^E(\mathbf D(\tau),V(\tau)),\\
D_{+-}'(\tau) &= -\beta D_{+-}(\tau)
    +H_+^E(\mathbf D(\tau),V(\tau)),\\
D_{-+}'(\tau) &= -\beta D_{-+}(\tau)
    +H_-^E(\mathbf D(\tau),V(\tau)),\\
D_{--}'(\tau) &= -\beta D_{--}(\tau)
    +H_-^E(\mathbf D(\tau),V(\tau)),\\
V'(\tau)
&=
-\kappa_zV(\tau)
+\frac{1}{2}\sigma_z^2V(\tau)^2
-s_k\bar\lambda\big(u_+^B+u_-^B\big),\\
C'(\tau)
&=
\kappa_z\bar zV(\tau)
+s_k\bar\lambda\Big(
H_+^B(\mathbf D(\tau),V(\tau))+u_+^B
+H_-^B(\mathbf D(\tau),V(\tau))+u_-^B
\Big)\\
&\quad
+\nu_k^{\mathrm{known}}\ell_{\mathrm{known}}
\left(q,\mathbf D(\tau)\cdot\mathbf r\right)
+\nu_k^{\mathrm{surprise}}
\left(
M_{\mathrm{surprise}}\left(q,\mathbf D(\tau)\cdot\mathbf r\right)-1
\right).
\end{aligned}
\]

\section{Derivation}

Use
\[
f_\tau(\mathbf y,z)
=
\exp\!\left(C(\tau)+\mathbf D(\tau)\cdot\mathbf y+V(\tau)z\right).
\]
The left side of the Kolmogorov equation gives
\[
\frac{\partial_\tau f_\tau}{f_\tau}
=
C'(\tau)+\mathbf D'(\tau)\cdot\mathbf y+V'(\tau)z.
\]

Hawkes decay contributes
\[
-\beta\mathbf D(\tau)\cdot\mathbf y.
\]

A buy event has multiplier
\[
z_+\exp(\mathbf D(\tau)\cdot\mathbf k_+ + V(\tau)\chi_+),
\]
and a sell event has multiplier
\[
z_-\exp(\mathbf D(\tau)\cdot\mathbf k_- + V(\tau)\chi_-).
\]
Subtracting one gives the usual jump-generator increment. Subtracting the
activity clock loading gives the Hamiltonian associated with the clock type.
For endogenous activity, this yields \(H_+^E\) multiplying
\(Y^{++}+Y^{+-}\) and \(H_-^E\) multiplying \(Y^{-+}+Y^{--}\). Combining with
Hawkes decay gives the four \(\mathbf D\) equations.

For baseline arrivals, the event intensities are both \(s_k\bar\lambda\).
Their pure event contribution is
\[
s_k\bar\lambda
\left[
z_+\exp(\mathbf D\cdot\mathbf k_+ + V\chi_+) -1
+z_-\exp(\mathbf D\cdot\mathbf k_- + V\chi_-) -1
\right].
\]
Because \(H_\pm^B+u_\pm^B=z_\pm\exp(\mathbf D\cdot\mathbf k_\pm+V\chi_\pm)-1\),
this is the baseline term in \(C'\):
\[
s_k\bar\lambda\left(H_+^B+u_+^B+H_-^B+u_-^B\right).
\]

The baseline return clocks are long-variance scaled, so their transform
loadings contribute
\[
-s_k\bar\lambda(u_+^B+u_-^B)z.
\]
This term is \(z\)-linear and therefore enters \(V'\), not \(C'\).

The continuous long-variance generator contributes
\[
\kappa_z\bar zV(\tau)
+\left(
-\kappa_zV(\tau)+\frac{1}{2}\sigma_z^2V(\tau)^2
\right)z.
\]
The constant part enters \(C'\), and the \(z\)-linear part enters \(V'\).
Adding the baseline clock term gives the stated Riccati equation for \(V\).

For news component \(c\), the activity jump and return jump have conditional
mark transform
\[
M_c\left(q,\mathbf D(\tau)\cdot\mathbf r\right)
=
\left(
1-\theta_U\left(
(\mathbf D(\tau)\cdot\mathbf r)a_c
+qm_c+\frac{1}{2}q^2\sigma_c^2
\right)
\right)^{-k_U}.
\]
Known-news timestamps are deterministic conditional on the schedule, so their
equivalent rate contributes the log transform
\[
\nu_k^{\mathrm{known}}\ell_{\mathrm{known}}
\left(q,\mathbf D(\tau)\cdot\mathbf r\right).
\]
Surprise news is compound Poisson, so it contributes
\[
\nu_k^{\mathrm{surprise}}
\left(
M_{\mathrm{surprise}}
\left(q,\mathbf D(\tau)\cdot\mathbf r\right)-1
\right).
\]
Both news terms are state independent after conditioning, so both enter
\(C'\).

Collecting the coefficients of \(1\), each Hawkes coordinate, and \(z\) gives
the stated system.

\section{Domain and validity}

The ODE system is valid on the maximal interval where \(C(\tau)\),
\(\mathbf D(\tau)\), and \(V(\tau)\) remain finite and the news mark
transforms stay inside the gamma-normal moment domain. The terminal conditions
must be finite. The model parameters must satisfy the Hawkes, news,
long-variance, and first-moment constraints listed in
\texttt{06\_affine\_domain\_constraints.tex}.

\section{Consequences}

\begin{itemize}
    \item The only nonlinear coefficient equation is the Riccati equation for \(V\), plus the exponential Hamiltonians in the Hawkes loading equations.
    \item Baseline event arrivals affect \(C\), but baseline return-clock exposure affects \(V\) because it is multiplied by \(Z_t\).
    \item Endogenous activity exposure affects \(\mathbf D\) because endogenous event intensity is state-linear.
    \item Known and surprise news differ only in the arrival transform, not in the mark transform itself.
\end{itemize}

\section{Checks}

\begin{tabularx}{\textwidth}{XXX}
\toprule
Check & Expected result & Evidence location\\
\midrule
Schedule RHS equations
& \(D'\), \(V'\), and \(C'\) are assembled with the Hamiltonians, long-variance drift, baseline clock loading, and news term shown here.
& \texttt{evaluate\_schedule\_base\_rhs} in \texttt{crates/janiform-affine-oracle/src/affine\_oracle/schedule/core.rs}\\
Event Hamiltonians with \(V\chi_\pm\)
& Buy and sell event multipliers include both \(\mathbf D\cdot\mathbf k_\pm\) and \(V\chi_\pm\).
& \texttt{evaluate\_long\_variance\_schedule\_kamiltonians} in \texttt{crates/janiform-affine-oracle/src/affine\_oracle/schedule/core.rs}\\
Known-news arrival transform
& Known news uses the deterministic-arrival log transform.
& Test \texttt{known\_news\_uses\_deterministic\_arrival\_log\_transform}\\
Surprise-news arrival transform
& Surprise news uses the compound-Poisson transform \(M-1\).
& Test \texttt{surprise\_news\_uses\_compound\_poisson\_transform}\\
Schedule tangent consistency
& Differentiating the same ODE system gives sensitivities matching finite differences.
& Test \texttt{schedule\_cumulant\_argument\_tangent\_matches\_finite\_difference}\\
\bottomrule
\end{tabularx}

\end{document}
